\section{问题重述与研究设计}
\subsection{任务与需要分开的时间对象}
2026年8月26日10:30，高位冰崩引发的碎屑流、泥石流与山洪链式过程冲击吉隆口岸。此类高位冰崩—碎屑流—泥石流链在喜马拉雅山区多次成灾，其运动性受冰水含量、体积与沿程地形共同影响\cite{schneider2011,iverson1997,richardson2000}。题给时间是口岸冲击锚点，并非源区冰体最初失稳的直接观测。研究需回答四个相连的问题：气象与近地层环境如何关联灾害，何种水源和床面条件达到触发门槛，何时可能触发及如何给出区间与预警，过程在既定地形和参数下如何演化。

令$t_m$为源区融水起始代理，$t_{\mathrm{trig}}$为真实泥石流形成时刻，$t_a^{(0)}$为同一成员在即时释放假设下传播到口岸代理的条件到达时刻，$t_a$为实际到达，$t_p$为固定发布时刻，$t_{\mathrm{hit}}$为题给口岸冲击锚点。冰体失稳可能发生在融水阈值之前，$t_m$与$t_{\mathrm{trig}}$并不等同。为避免同一次事件参与自身验证，本文固定$t_p=$2026-08-26 08:00当地时，先产生源区与口岸配对成员，再以10:30作事后单点核对。四个问题分别回答环境异常、力平衡状态、条件时间区间和传播情景，辅助包含只作单点核对。

地形、物源和水源分别对应能量、物质与触发条件。本文以水源供给与床面起动比值作为主状态，避免把资料条件不同的过程压缩为单一概率；空间模型用于环境对照，运动模型给出条件传播，各问共用同一条公共中心线。

\subsection{资料、来源与用途}
题给八类资料、其发布机构与本文用途见表\ref{tab:data}。来源以机构名与产品名注明，请求参数与产品版本随成果保存。

\begin{table}[htbp]\centering\footnotesize
\caption{题给资料及用途}\label{tab:data}
\begin{tabularx}{\textwidth}{Y >{\hsize=1.10\hsize}Y >{\hsize=1.05\hsize}Y >{\hsize=0.85\hsize}Y}\toprule
附件&资料范围&来源产品&用途与限制\\\midrule
1&2021--2026逐时MET/RADM/SOIL/FLUX&吉隆灾害专题（国家青藏高原科学数据中心）&温湿风廓线、辐射、土壤及湍流诊断；单站代表性有限\\
2&2005--2025周边县日均温（至2025-08）&ERA5-Land逐日统计产品，同上&长期气候背景和同口径积温对照，不代替源区温度观测\\
3&2000--2020四期高程变化&全球冰川高程变化产品，同上&有效像元加权区域背景，均值误差不视为区域均值标准误\\
4&灾前冰川DEM、冰厚、速度&冰川地形与冰厚遥感产品，同上&原生源区积分、源几何与速度重叠旁证\\
5&2020--2023土地覆盖&土地覆被遥感产品，同上&原类码作为分类变量；未确认类码不解释为建设用地\\
6&84处历史灾点面及属性&历史灾点调查数据，同上&沟道临界域、规模排序、历史等效摩擦和面内像元对照\\
7&事件日期研究域DEM&事件区数字高程模型，同上&公共水文地形与路由载体；日期不证明灾前或灾后采集\\
8&第七次县级人口普查&县级普查汇总，同上&县级承灾背景，不空间化为精确暴露人口\\\bottomrule
\end{tabularx}\end{table}

融水面积权重取附件3的RGI6编目多边形与集水区的面积交比，属于编目年代边界，不是2026年实时冰川或积雪范围；高程带代理并列报告但不进入融水。补充的ERA5与MERRA2用于同一UTC+8自然日的产品对照；IFS采用按单次起报保存的06、12、18 UTC档案，分别记录起报时刻、可用上界、有效小时和产品版本。起报可用性表示信息时次满足发布要求，源区冰面实况仍需独立观测确认。

站表共49545行，物理有效46755小时；五层温湿共同有效45577小时，可计算整体理查森数的共同样本44018小时。历史84处灾点中降雨型41、冰水混合型25、冰雪融水型16、冰湖溃决型2，大及特大规模合计9处。该目录每个对象都已发生过灾害，因此无法据此构造完整的发生/未发生样本；DEM相交灾点只是其中一部分，像元样本量不等于独立灾害数量。人口“合计”字段为17536人，只给县级背景，不按均匀密度分摊到狭窄演算走廊。

\subsection{坐标与位置的来源说明}
题面只给“吉隆口岸”，未给实际受灾断面坐标。本文以(85.3826$^\circ$E,28.2826$^\circ$N)为工作点，以(85.3777$^\circ$E,28.2797$^\circ$N)为参考候选，两者平面相距约571米，均按DEM近似位置处理。固定源区、网格、体积和动力参数，仅改变目标点重演：两候选的路线吸附偏移为5.186582和27.40米，D8到达为32.766603和33.90分钟，峰流深10.085302和7.42米，峰速6.601474和8.82米每秒。5.19米是路线吸附偏移，不代表口岸定位精度；两候选间距另行报告。

站点坐标(85.30$^\circ$E, 28.40$^\circ$N)是吉隆沟高寒山区的假定代表点：原始MET、RADM、SOIL与FLUX档案均不含站点元数据。海拔取气压反演值3978.85\,m（标准大气推断），不是实测安装高度；辐射与触发计算取整使用3980\,m，与反演值相差1.15\,m，远小于反演不确定度。太阳能时角计算使用该经度，经度假定误差会线性传播至辐射诊断。

\subsection{质量控制和空间一致化}
四类站点资料先按时间戳拼接、去重并重建连续小时轴；哨兵值及不合理物理值置空，质量码释义未充分提供时保守屏蔽非零码。滚动三小时要求三小时均有效，三日小时正积温要求72小时完整；普通自然日聚合至少20个有效小时。日光反照率和晴空指数只在符合日光条件的时段考察，采用有效比例要求而非要求不存在的20小时日照。无观测小时不补零降水或融化；跨缺测的API先重启并保留暖启动期，雨事件分割只把完整观测的干期当间隔，不完整过程的水源状态记为未知。该处理减少可用样本，但避免把缺测误作无雨无融；插补、分位与标准化只在训练集拟合，事件日全日资料仅作描述，不作为事前特征。

空间资料统一以像元中心定位。经纬网格的米制尺寸近似为
\begin{equation}
 d_y=110574|\Delta\phi|,\qquad
 d_x=111320\cos\phi_0|\Delta\lambda|,\quad \phi_0=28.4^\circ.
\end{equation}
主网格南北约30.71米、东西约27.20米，面积约835.27平方米，与900平方米的名义值有明显差异。冰川原生30米UTM网格上的体积先积分再投影，防止多个原生像元落到同一演算格后重复积分；粗化网格原点位于整个块的中心而非左上第一个细格中心。D8选向按真实链接长度比较最大坡降，与后续流量、时间和CFL所用几何一致。

\subsection{外部气象对照}
请求明确指定ERA5、UTC与关闭海拔自动降尺度；NASA POWER请求明确指定UTC。两套逐时响应先转为UTC+8，再按完整24小时汇总事件自然日。
\begin{table}[htbp]\centering\small
\caption{同日界外部气象对照}\label{tab:extmet}
\begin{tabular}{@{}l c c c@{}}\toprule
产品&响应代表高程/m&日均/最高/最低温度/$^\circ$C&日降水/mm\\\midrule
ERA5（Open-Meteo接口）&4322&9.41/13.10/6.40&7.50\\
MERRA2（NASA POWER接口）&4836.18&6.09/11.06/2.43&1.22\\\bottomrule
\end{tabular}\end{table}
两产品事件日降水差异明显，日均温相差约3.32 K。格点代表高程相差约514 m，但温度差与高程差不同向，无法据此反推直减率；粗格点、下垫面与模式偏差未分离，两个接口也不保证统计独立，因此只作并列对照，不替代站点观测。

第三类资料是IFS单次起报。主起报为2026年8月25日18 UTC，并以同日06、12 UTC作模式演化对照；按六小时分发上界，三者最迟分别于25日12 UTC、18 UTC和26日00 UTC可用，均不晚于固定发布时刻。每次取发布前最后完成小时和其后24个有效小时，源区与口岸配对生成1000个成员。普通短波与降水按原始结束标签对齐到$[t,t+1)$，温湿取两端时点值的梯形均值；任一端点缺失则该小时不可评估，不用下一条记录跨过时间缺口。18 UTC口岸格点序列只作独立上下文，不进入主传播链。该流程保证信息时次不晚于发布；模式输出与冰面实况的一致性需另行检验。

站点对照按UTC+8转UTC日，仅保留温降水均有24个有效小时的共同日期；2021--2025分别保留247、363、361、363、214日。2025同窗口站点降水102.00毫米、POWER为742.41毫米；站点年总量不完整，二者不作同口径比较。站点位置、时区与小时标签仍是假设，故只是代表点产品差异。

原始响应、请求参数与产品版本随成果保存；公开站点元数据缺失，真实安装高度及时区仍未知，海拔按反演值处理，不视为实测信息。

\section{问题分析与建模思路}
四个问题共享同一套资料、时间轴和公共几何，但回答对象不同：问题一识别事件前环境异常，问题二判断水源与床面是否达到起动条件，问题三在固定发布时刻给出条件时间产品，问题四在给定体积与地形下演算传播。图\ref{fig:roadmap}给出全文技术路线。

四问共享时间基准、坐标和注入断面，箭头只表示计算依赖；本文将结果分为计算闭合、模型内的条件分布与真实灾害外推三层，结论限于前两层。

\begin{figure}[H]\centering
\begin{tikzpicture}[
  font=\footnotesize,
  base/.style={draw=black!55, rounded corners=2pt, align=center, inner sep=3pt},
  data/.style={base, fill=gray!8, text width=13.8cm, minimum height=1.05cm},
  q/.style={base, fill=teal!7, text width=3.30cm, minimum height=3.55cm, font=\scriptsize},
  outbox/.style={base, fill=orange!8, text width=13.8cm, minimum height=1.2cm},
  arr/.style={-{Latex[length=2mm]}, semithick, draw=black!65}
]
\node[data] (data) at (0,0) {\textbf{题给资料与公共预处理}\\四类逐时观测、DEM、冰川编目、历史灾点、IFS\\时间基准对齐｜质量控制｜公共注入断面｜源区体积};
\node[q] (q1) at (-5.1,-2.65) {\textbf{问题一}\\环境关联\\31项候选\\稳健异常→相关簇\\正式前五};
\node[q] (q2) at (-1.7,-2.65) {\textbf{问题二}\\起动判别\\床面受力→水控/\\重力分支\\事件累计SCS-CN\\→$T(t)$};
\node[q] (q3) at (1.7,-2.65) {\textbf{问题三}\\条件预报\\固定08:00\\3000成员\\$t_m$、$t_a^{(0)}$区间};
\node[q] (q4) at (5.1,-2.65) {\textbf{问题四}\\情景传播\\一维自然停止\\D8路由\\首达、峰值、右删失};
\node[outbox] (ev) at (0,-5.75) {\textbf{结果层次与适用范围}\\计算闭合（单位/状态/账本/CFL）→ 条件成员分布 → 真实触发外推\\真实 $t_r$、$t_f$、$t_a$、真实预警与覆盖率尚未给出};
\draw[arr] (data.south) -- (q1.north);
\draw[arr] (data.south) -- (q2.north);
\draw[arr] (data.south) -- (q3.north);
\draw[arr] (data.south) -- (q4.north);
\draw[arr] (q1.south) -- (ev.north -| q1.south);
\draw[arr] (q2.south) -- (ev.north -| q2.south);
\draw[arr] (q3.south) -- (ev.north -| q3.south);
\draw[arr] (q4.south) -- (ev.north -| q4.south);
\end{tikzpicture}
\caption{全文建模思路与结果层次}\label{fig:roadmap}
\end{figure}


\section{基本假设与符号约定}
\subsection{公共假设与适用范围}
跨问题主要符号见表\ref{tab:symbols}；公共假设逐条列于其后，编号H0-1至H0-9。

\begin{table}[htbp]\centering\footnotesize
\caption{主要符号}\label{tab:symbols}
\begin{tabular}{@{}l l l@{}}\toprule
符号&意义&单位\\\midrule
$\theta_v,Ri_b$&虚位温、整体理查森数&K、无量纲\\
$R_n,H,LE,G$&净辐射、感热、潜热、土壤热通量&W/m$^2$\\
$e_m,m_{src},m_{cat}$&小时能量、源区/流域等效融水&MJ/m$^2$、mm\\
$Q_w,Q_c,q_c$&水源流量、起动流量、单宽起动流量&m$^3$/s、m$^2$/s\\
$T,T_w,P$&触发指数、双通道候选指示、单调相对评分&无量纲\\
$S,n,\tau_y,\mu_H$&比降、糙率、屈服门槛、等效摩擦&1、s/m$^{1/3}$、Pa、1\\
$V_{src},V_{rel},V_{df}$&源冰总量、释放冰量、混合入沟体积&m$^3$\\
$h,v,I$&等效泥深、速度、单位宽动量通量&m、m/s、N/m\\
$t_m,t_{\mathrm{trig}},t_a^{(0)},t_a$&融水起始、真实形成、条件到达、实际到达&h\\
$t_p,t_{\mathrm{hit}}$&固定发布、题给冲击锚点&h\\
$V_{inj},V_{dom},V_{out}$&累计注入、域内、累计出域体积&m$^3$\\
$\mu_{dyn},\kappa,v_0;\ \Delta V_{a,j},B_{a,j}$&动力与节点库存&1、m$^{-1}$、m/s、m$^3$\\\bottomrule
\end{tabular}\end{table}

\begin{enumerate}
\item 题给冲击时刻与附件逐时记录按同一UTC+8标度对齐；IFS原始有效时刻按UTC保存，再转换到当地发布窗。时区元数据缺失，时间对齐仍是主要不确定性。
\item 单站代表同一高程带的大尺度强迫，源区气温按0.5--0.7 K/(100 m)直减率做外推敏感性；不假定局地风和冰体温度已知。
\item 能量余量表示潜在融化能量，闭合缺口不全部归因于实际消融。冰厚、源区比例与物源放大均作为明示场景参数。
\item 模型成员概率、条件分位、确定配置包络和真实事件频率分别命名，避免混合成单一的真实灾害置信区间。
\item 前向链的$t_m$是融水起始代理，$t_{\mathrm{trig}}$是真实泥石流形成时刻；$t_a^{(0)}$由同一成员经固定传播配置得到，是即时释放假设下的条件到达，$t_a$是实际到达。真实释放时刻$t_r$、形成时刻$t_{\mathrm{trig}}$和实际到达$t_a$缺少独立观测，本文暂不给出估计。
\item 过程引擎采用固定地形和集中有限历时注入；侵蚀补给、颗粒分选和堰塞湖溃坝未显式求解。D8终止于分析域时记为右删失；已有漫溢背景作为后续外加源风险，不计入当前释放边界。
\item 历史沟长和扇长按水平投影解释；DEM语义不明时并列冰面与冰下高程假设，不用高分位消除语义歧义。
\item 有限库存对照中的沿程节点与库存均为合成机制情景，不代表已确认的物源点；其速度是团块沿固定路线的条件速度，不等于真实前锋速度或流深。
\item 口岸工作坐标为DEM量图近似位置（28.2826$^\circ$N, 85.3826$^\circ$E），与参考候选相差约0.6\,km；源区至注入断面水平距离2041.449\,m，公共注入至口岸代理中心线20882.237\,m。站点坐标为假定代表点，海拔取气压反演值3978.85\,m。上述位置均按DEM近似，不是实测安装或受灾断面数据。
\end{enumerate}

这些假设由现有资料条件决定；若与观测冲突，本文保留原设定并报告偏差，不通过调整起点、剔除成员或放宽阈值追求表面一致。阈值用于比较，进一步率定仍需独立灾害目录和现场观测。

全文小时步上的mm是该小时水当量深，不额外再除一小时；流量转换统一使用$1\mymunit{mm}\times1\mymunit{km^2}/1\mymunit{h}=1/3.6\mymunit{m^3/s}$。同一符号在日累计时附日下标，避免把全期累计能量与单日门槛相比较。
